\documentclass{article}

% NeurIPS / ICML / ICLR Style Formatting Setup
\usepackage[utf8]{utf8}
\usepackage[margin=1in]{geometry}
\usepackage{amsmath,amssymb,amsfonts,amsthm}
\usepackage{graphicx}
\usepackage{booktabs}
\usepackage{hyperref}
\usepackage{cite}
\usepackage{microtype}
\usepackage{xcolor}

\title{\textbf{Optimization Conflicts in Continual Multi-Task Learning: \\ Why Replay Degrades Regularization-Based Knowledge Retention}}

\author{
  \textbf{CDP Research Team} \\
  Scalable Customer Data Platform Lab \\
  \texttt{research@cdp-platform.io}
}

\date{\today}

\begin{document}

\maketitle

\begin{abstract}
Continual learning in deployed multi-task architectures requires balancing knowledge preservation on historical task distributions with plastic adaptation to emergent domain shifts. While parameter regularization (e.g., Elastic Weight Consolidation, EWC) and experience replay buffer sampling are individually established paradigms for mitigating catastrophic forgetting, their interaction in multi-task neural networks remains under-explored. In this work, we investigate the continual deployment of a Multi-Task Learning (MTL) neural network predicting customer churn (binary classification) and Customer Propensity Index (CPI, regression) across non-stationary customer contract segments. Through an empirical ablation study ($N=3,875$ initial task samples, $\lambda_{\text{ewc}}=100.0$, buffer size $M=1,000$, replay ratio $r=0.20$), we uncover a striking pathology: \textbf{experience replay alone severely exacerbates catastrophic forgetting (+10.11\% forgetting rate vs +2.41\% baseline), while combining EWC with Replay degrades retention performance compared to EWC alone (-1.41\% vs -2.44\% forgetting rate) and harms target task plasticity ($\text{AUC}_B$ drops from 0.8162 to 0.8005).} We formalize four investigative hypotheses explaining this optimization conflict---Gradient Interference, Over-Regularization under Stability-Plasticity Strain, Taylor Basins Drift, and Buffer Overfitting---providing a theoretical foundation and practical diagnostic guidelines for continual multi-task deployment.
\end{abstract}

\section{Introduction}
Industrial predictive platforms increasingly rely on Multi-Task Learning (MTL) deep neural architectures to simultaneously optimize interrelated business objectives, such as binary churn classification and continuous engagement or value modeling \cite{caruana1997multitask}. When deployed in real-world environments, these systems encounter distribution shifts over time as customer demographics, contract structures, or economic conditions evolve. Naive fine-tuning on sequential domains induces \textbf{catastrophic forgetting}, wherein gradient updates for a new task $D_B$ overwrite parameter representations critical for historical task $D_A$ \cite{french1999catastrophic, mccloskey1989catastrophic}.

To mitigate catastrophic forgetting, two predominant continual learning paradigms have emerged:
\begin{enumerate}
    \item \textbf{Regularization-Based Approaches:} Methods such as Elastic Weight Consolidation (EWC) penalize updates to parameters proportional to their historical importance, estimated via the diagonal of the Fisher Information Matrix \cite{kirkpatrick2017overcoming}.
    \item \textbf{Replay-Based Approaches:} Experience replay preserves a subset of historical exemplars in a memory buffer $M$ and interleaves them into batches of incoming task data \cite{chaudhry2019efficient, riemer2019learning}.
\end{enumerate}

While prior literature frequently assumes that parameter regularization and experience replay are complementary mechanisms that can be naively combined \cite{nguyen2018variational}, their joint behavior within \emph{multi-task architectures} under non-stationary task transitions remains under-investigated. Multi-task loss formulations introduce shared representation dynamics where shared backbone parameters $\theta_{\text{shared}}$ are simultaneously constrained by multiple task heads ($f_A$ and $f_B$).

\paragraph{Core Research Question:}
\emph{How do parameter-regularization penalties (EWC) and experience replay interact when applied to multi-task neural architectures subject to non-stationary task segment shifts? Specifically, does experience replay enhance or interfere with EWC's ability to preserve knowledge across shared representations?}

\subsection{Key Contributions}
\begin{enumerate}
    \item \textbf{Empirical Discovery of Optimization Conflict:} We establish that combining EWC ($\lambda=100.0$) and stratified experience replay ($M=1000, r=0.20$) in an MTL architecture produces sub-optimal retention compared to EWC alone, degrading Task A retention from $-2.44\%$ (EWC alone) to $-1.41\%$ (Full EWC+Replay) and reducing Task B plasticity ($\text{AUC}_B$ 0.8162 vs 0.8005).
    \item \textbf{Pathology of Pure Replay in MTL:} We show that experience replay \emph{without} parameter regularization causes severe catastrophic forgetting ($+10.11\%$ forgetting rate), performing significantly worse than naive fine-tuning ($+2.41\%$), demonstrating that small replay buffers can act as non-stationary noise injections in MTL loss surfaces.
    \item \textbf{Four Theoretical Hypotheses:} We construct a formal framework detailing four potential mechanisms for this optimization conflict: (i) \emph{Gradient Interference \& Optimization Noise}, (ii) \emph{Stability-Plasticity Over-Regularization}, (iii) \emph{Taylor Basins Drift}, and (iv) \emph{Buffer Overfitting}.
    \item \textbf{Architectural Guidelines for Continual MTL:} We provide actionable recommendations for machine learning practitioners deploying MTL pipelines in production, cautioning against naive hybridization of EWC and experience replay.
\end{enumerate}

\section{Related Work}

\subsection{Multi-Task Learning (MTL) and Loss Weighting}
Multi-Task Learning aims to improve generalization by leveraging domain-specific information contained in the training signals of related tasks \cite{caruana1997multitask, ruder2017overview}. In deep neural networks, hard parameter sharing exposes a joint parameter backbone $\theta_{\text{shared}}$ to gradients originating from distinct task heads. A foundational challenge in MTL is handling conflicting task gradients, where updates for Task 1 oppose updates for Task 2 ($\langle \nabla \mathcal{L}_1, \nabla \mathcal{L}_2 \rangle < 0$), leading to sub-optimal Pareto multi-task solutions \cite{yu2020gradient, sener2018multi}.

\subsection{Continual Learning \& Catastrophic Forgetting}
Continual learning algorithms operate on sequential task streams $D_1, D_2, \dots, D_T$ without storing full historical datasets. Continual learning approaches broadly fall into three categories:
\begin{itemize}
    \item \textbf{Regularization Methods:} EWC \cite{kirkpatrick2017overcoming}, Synaptic Intelligence \cite{zenke2017continual}, and Memory Aware Synapses \cite{aljundi2018memory} penalize parameter displacement along directions sensitive to prior tasks.
    \item \textbf{Replay Methods:} Experience Replay \cite{chaudhry2019efficient} and Gradient Episodic Memory (GEM / A-GEM) \cite{lopez2017gradient, chaudhry2018efficient} store historical exemplars or project gradient updates to satisfy historical task constraints.
    \item \textbf{Architecture-Based Methods:} Progressive Neural Networks \cite{rusu2016progressive} allocate dedicated subnetworks per task.
\end{itemize}

\subsection{Optimization Conflicts in Hybrid Continual Learning}
While hybrid approaches combining EWC and Replay have been explored in single-task vision benchmarks \cite{nguyen2018variational, riemer2019learning}, their interactions in MTL networks remain largely unexamined. Riemer et al. \cite{riemer2019learning} demonstrated that experience replay can induce gradient alignment issues if replay vectors conflict with current task gradients. However, the theoretical interaction between quadratic EWC parameter constraints and loss gradients derived from small multi-task replay buffers under severe domain shifts has not been formally characterized in industrial tabular applications.

\section{Methodology \& Theoretical Framework}

\subsection{Task Sequence and MTL Architecture}
We formulate customer domain progression as a continual learning stream over sequential contract segments:
\begin{itemize}
    \item \textbf{Task A ($D_A$):} Month-to-month contracts ($N = 3,875$ samples), characterized by high churn variance and high churn baseline.
    \item \textbf{Task B ($D_B$):} Long-term contracts (1-year and 2-year contracts), featuring a low churn baseline and a shifted feature distribution.
\end{itemize}

The model is a deep Multi-Task Learning neural network parameterized by $\theta = \{\theta_{\text{shared}}, \theta_A, \theta_B\}$, comprising:
\begin{itemize}
    \item A shared feature extractor backbone: $h = g(X; \theta_{\text{shared}})$
    \item Head A (Binary Churn Classification): $f_A(h; \theta_A) \to \hat{y}_{\text{churn}} \in \mathbb{R}$
    \item Head B (Customer Propensity Index, CPI Regression): $f_B(h; \theta_B) \to \hat{y}_{\text{cpi}} \in \mathbb{R}$
\end{itemize}

The joint multi-task loss on a mini-batch $B$ is defined as:
\begin{equation}
\mathcal{L}_{\text{MTL}}(\theta; B) = \alpha \cdot \mathcal{L}_{\text{BCE}}\left(f_A(X; \theta), y_{\text{churn}}\right) + \beta \cdot \mathcal{L}_{\text{MSE}}\left(f_B(X; \theta), y_{\text{cpi}}\right)
\end{equation}
where $\alpha = 0.7$ and $\beta = 0.3$ represent static task loss weights.

\subsection{Continual Learning Formulations}

\subsubsection{1. Baseline (Fine-Tuning)}
When transitioning from Task A to Task B, naive fine-tuning optimizes:
\begin{equation}
\theta_B^* = \arg\min_{\theta} \mathcal{L}_{\text{MTL}}(\theta; D_B)
\end{equation}
initialized at optimal parameters $\theta_A^* = \arg\min_{\theta} \mathcal{L}_{\text{MTL}}(\theta; D_A)$.

\subsubsection{2. Elastic Weight Consolidation (EWC)}
EWC adds a quadratic parameter penalty centered at $\theta_A^*$:
\begin{equation}
\mathcal{L}_{\text{EWC}}(\theta; D_B) = \mathcal{L}_{\text{MTL}}(\theta; D_B) + \frac{\lambda_{\text{ewc}}}{2} \sum_{i} F_{A,i} (\theta_i - \theta_{A,i}^*)^2
\end{equation}
where $\lambda_{\text{ewc}} = 100.0$, and $F_{A,i}$ is the $i$-th diagonal element of the Fisher Information Matrix calculated over Task A:
\begin{equation}
F_{A,i} = \frac{1}{|D_A|} \sum_{x \in D_A} \left( \frac{\partial \mathcal{L}_{\text{MTL}}(x; \theta_A^*)}{\partial \theta_i} \right)^2
\end{equation}

\subsubsection{3. Experience Replay}
A memory buffer $M_A$ of size $M = 1000$ samples is populated from $D_A$ using stratified sampling across churn labels. During training on Task B, each mini-batch $B$ is constructed by mixing $(1-r)$ fraction of current Task B data with $r = 0.20$ fraction of replay data sampled from $M_A$:
\begin{equation}
B_{\text{mixed}} = (1-r) B_B \cup r B_{\text{replay}}, \quad \mathcal{L}_{\text{Replay}}(\theta) = \mathcal{L}_{\text{MTL}}(\theta; B_{\text{mixed}})
\end{equation}

\subsubsection{4. Full Hybrid (EWC + Replay)}
The combined objective minimizes EWC parameter regularization alongside mixed replay batch optimization:
\begin{equation}
\mathcal{L}_{\text{Full}}(\theta) = \mathcal{L}_{\text{MTL}}(\theta; B_{\text{mixed}}) + \frac{\lambda_{\text{ewc}}}{2} \sum_{i} F_{A,i} (\theta_i - \theta_{A,i}^*)^2
\end{equation}

\subsection{Investigative Hypotheses for Optimization Conflict}

\paragraph{Hypothesis 1: Gradient Conflict and Optimization Noise ($\mathcal{H}_1$).}
When optimizing $\mathcal{L}_{\text{Full}}(\theta)$, the parameter update vector decomposes into three force components:
\begin{equation}
\nabla \mathcal{L}_{\text{Full}}(\theta) = (1-r) \nabla \mathcal{L}_{\text{MTL}}(\theta; B_B) + r \nabla \mathcal{L}_{\text{MTL}}(\theta; B_{\text{replay}}) + \lambda_{\text{ewc}} F_A \odot (\theta - \theta_A^*)
\end{equation}
We hypothesize that the replay gradient $\nabla \mathcal{L}_{\text{replay}}$ and the EWC quadratic pull $\nabla \mathcal{R}_{\text{EWC}}$ exhibit negative directional cosine similarity:
\begin{equation}
\cos\left( \nabla \mathcal{L}_{\text{replay}}, \nabla \mathcal{R}_{\text{EWC}} \right) < 0
\end{equation}
This gradient misalignment causes high-frequency stochastic oscillation along the shared manifold $\theta_{\text{shared}}$, creating optimization noise that counteracts EWC's smooth quadratic constraint.

\paragraph{Hypothesis 2: Stability-Plasticity Strain \& Over-Regularization ($\mathcal{H}_2$).}
A large penalty coefficient ($\lambda_{\text{ewc}} = 100.0$) severely constrains the feasible parameter update region $\mathcal{S}_{\epsilon} = \{\theta \mid \frac{1}{2} (\theta - \theta_A^*)^T F_A (\theta - \theta_A^*) \le \epsilon\}$. Forcing the model to simultaneously satisfy Task B objective $\mathcal{L}(D_B)$ and buffer objective $\mathcal{L}(M_A)$ inside this tightly bounded set over-regularizes the network, leading to suboptimal local minima that compromise both historical retention and new task plasticity.

\paragraph{Hypothesis 3: Breakdown of First-Order Taylor Approximation / Taylor Basins Drift ($\mathcal{H}_3$).}
EWC relies on a second-order Taylor expansion around $\theta_A^*$:
\begin{equation}
\mathcal{L}_A(\theta) \approx \mathcal{L}_A(\theta_A^*) + \frac{1}{2} (\theta - \theta_A^*)^T F_A (\theta - \theta_A^*)
\end{equation}
As gradient updates driven by $B_{\text{mixed}}$ displace parameters away from $\theta_A^*$, higher-order residual terms $O(\|\theta - \theta_A^*\|^3)$ become dominant. Consequently, $F_A$ calculated at $\theta_A^*$ ceases to represent true curvature in the displaced region, causing EWC penalty calculations to misguide parameter optimization.

\paragraph{Hypothesis 4: Buffer Overfitting and Distributional Bias ($\mathcal{H}_4$).}
A replay buffer size of $M = 1000$ represents only $25.8\%$ of $D_A$ ($N = 3875$). Repeated optimization over this small subsample introduces buffer overfitting. The model fits specific exemplars in $M_A$ rather than the broader distribution $P(X_A, Y_A)$, resulting in decision boundary distortion on the complete Task A validation set.

\section{Empirical Evaluation \& Ablation Study}

\subsection{Experimental Setup}
\begin{itemize}
    \item \textbf{Dataset:} $N_A = 3,875$ Month-to-month contract samples (Task A) and Task B Long-term contract samples.
    \item \textbf{Optimization:} Adam optimizer, learning rate $\eta = 10^{-3}$, batch size $|B| = 64$, training epochs $E = 50$.
    \item \textbf{CL Setup:} EWC coefficient $\lambda_{\text{ewc}} = 100.0$, buffer size $M = 1000$ (stratified), replay ratio $r = 0.20$.
    \item \textbf{Forgetting Rate Metric (\%):} Defined as:
    \begin{equation}
    \text{Forgetting Rate (\%)} = \frac{\text{AUC}_{A, \text{before}} - \text{AUC}_{A, \text{after}}}{\text{AUC}_{A, \text{before}}} \times 100\%
    \end{equation}
\end{itemize}

\subsection{Core Empirical Results (Seed 42 Benchmark)}

\begin{table}[h]
\centering
\caption{\textbf{Empirical Ablation Study (Seed 42 Benchmark).} Evaluation of continual learning methods on the Multi-Task Churn and CPI model transitioning from Task A to Task B.}
\label{tab:ablation_results}
\vspace{0.2cm}
\begin{tabular}{lcccccc}
\toprule
\textbf{Configuration} & \textbf{Use EWC} & \textbf{Use Replay} & $\mathbf{\text{AUC}_{A, \text{before}}}$ & $\mathbf{\text{AUC}_{A, \text{after}}}$ & $\mathbf{\text{AUC}_{B}}$ & \textbf{Forgetting Rate (\%)} \\
\midrule
\textbf{Baseline (Fine-tune)} & ❌ & ❌ & 0.6830 & 0.6665 & \textbf{0.8463} & +2.41\% \\
\textbf{Replay Only} & ❌ & ✅ & 0.6830 & 0.6140 & 0.8246 & +10.11\% \\
\textbf{EWC Only} & ✅ & ❌ & 0.6830 & \textbf{0.6997} & 0.8162 & \textbf{-2.44\%} \\
\textbf{Full (EWC + Replay)} & ✅ & ✅ & 0.6830 & 0.6927 & 0.8005 & \textbf{-1.41\%} \\
\bottomrule
\end{tabular}
\end{table}

\subsection{Detailed Empirical Analysis}

\subsubsection{1. Failure of Replay Alone (+10.11\% Forgetting Rate)}
Replay alone degrades Task A performance significantly, dropping $\text{AUC}_A$ from 0.6830 to 0.6140 (+10.11\% forgetting rate). This is dramatically worse than naive fine-tuning (+2.41\% forgetting rate). In a multi-task tabular setting, mixing 20\% replay data without parameter constraints introduces non-stationary mini-batch variances. The shared backbone $\theta_{\text{shared}}$ receives competing signals from $B_B$ and $B_{\text{replay}}$, corrupting historical decision boundaries. This directly validates \textbf{Hypothesis 4 (Buffer Overfitting)} and \textbf{Hypothesis 1 (Gradient Interference)}.

\subsubsection{2. Superiority of EWC Alone (-2.44\% Forgetting Rate)}
EWC alone achieves the highest retention ($\text{AUC}_A$ increases from 0.6830 to 0.6997, yielding $-2.44\%$ forgetting rate). This negative forgetting rate demonstrates \textbf{positive backward transfer}, where training on Task B under quadratic Fisher constraints refines shared representations in a way that benefits Task A performance.

\subsubsection{3. Sub-Optimality of Hybrid EWC + Replay (-1.41\% vs -2.44\%)}
Combining EWC and Replay yields $\text{AUC}_{A,\text{after}} = 0.6927$ (-1.41\% forgetting rate), which is \textbf{1.03 percentage points worse in retention} than EWC alone. Furthermore, Task B plasticity is degraded to $\text{AUC}_B = 0.8005$, compared to 0.8162 for EWC alone and 0.8463 for Fine-tuning. This confirms that adding replay to an EWC-regularized multi-task model degrades both retention and plasticity, empirically validating the optimization conflict hypothesis!

\section{Discussion \& Practical Guidelines}

\subsection{Deconstructing the Optimization Conflict}
The empirical findings confirm that parameter-space regularization and data-space replay do not combine additively in multi-task networks. Instead, they exhibit negative interference:
\begin{enumerate}
    \item \textbf{Gradient Vector Misalignment:} EWC pulls parameters toward the historical Taylor basin $\theta_A^*$, while replay mini-batches push parameters toward current exemplars. As predicted by $\mathcal{H}_1$, these vector forces cancel out along critical gradient dimensions, leading to ineffective parameter updates.
    \item \textbf{Plasticity Constraining:} As predicted by $\mathcal{H}_2$, the combination of a tight quadratic penalty ($\lambda=100.0$) and a dual-task mixed batch severely restricts backbone flexibility, lowering Task B AUC from 0.8463 (Baseline) down to 0.8005 (Full).
\end{enumerate}

\subsection{Deployment Guidelines for Multi-Task Systems}
For machine learning engineers deploying continual MTL systems in production:
\begin{itemize}
    \item \textbf{Avoid Naive Hybridization:} Do not blindly combine EWC penalties with replay buffers without monitoring gradient alignment ($\cos(\nabla \mathcal{L}_{\text{replay}}, \nabla \mathcal{R}_{\text{EWC}})$).
    \item \textbf{Prefer EWC for Feature-Preserving Shifts:} When distribution shifts occur across structured domain segments (e.g., contract types), parameter regularization via Fisher importance provides superior retention and positive transfer.
    \item \textbf{Buffer Quality over Sample Quantity:} If experience replay must be used, sample buffer selection should prioritize boundary exemplars (e.g., via A-GEM or maximal interference sampling) rather than standard stratified sampling.
\end{itemize}

\section{Limitations and Future Work}
\subsection{Limitations}
\begin{enumerate}
    \item \textbf{Single Task Transition:} Empirical validation currently focuses on the transition $D_A \to D_B$. Long-horizon task chains ($D_A \to D_B \to D_C \dots$) require further investigation.
    \item \textbf{Diagonal Fisher Assumption:} EWC employs a diagonal approximation of the Fisher Information Matrix, ignoring off-diagonal parameter covariances.
    \item \textbf{Fixed Hyperparameters:} Experiments used fixed $\lambda_{\text{ewc}} = 100.0$ and replay ratio $r = 0.20$. Dynamic scheduling of these parameters remains an open direction.
\end{enumerate}

\subsection{Pre-Experimental \& Investigative Roadmap}
To further isolate the exact physical mechanism among our four hypotheses, upcoming experimental iterations will integrate:
\begin{itemize}
    \item \textbf{Gradient Cosine Tracking:} Continuously logging $\cos(\nabla \mathcal{L}_{\text{replay}}, \nabla \mathcal{R}_{\text{EWC}})$ across training epochs to directly test $\mathcal{H}_1$.
    \item \textbf{Taylor Loss Residual Analysis:} Measuring $O(\|\theta - \theta_A^*\|^3)$ drift to quantify $\mathcal{H}_3$.
    \item \textbf{Buffer Size Scaling Sweeps:} Evaluating $M \in \{250, 500, 1000, 2000, 3875\}$ to determine the exact threshold where buffer overfitting ($\mathcal{H}_4$) transitions to true distribution matching.
\end{itemize}

\section{Conclusion}
This paper presents an investigation into continual learning within multi-task customer churn architectures. Through an empirical ablation study across sequential customer contract segments, we demonstrate that combining Elastic Weight Consolidation (EWC) and experience replay leads to an unexpected optimization conflict. Experience replay alone causes severe catastrophic forgetting (+10.11\%), while adding replay to EWC degrades retention performance (-1.41\% vs -2.44\%) and reduces new task plasticity ($\text{AUC}_B$ 0.8005 vs 0.8162). We formalize four investigative hypotheses explaining this phenomenon and provide practical recommendations for robust continual MTL deployment.

\bibliographystyle{plain}
\begin{thebibliography}{10}

\bibitem{caruana1997multitask}
R.~Caruana.
\newblock Multitask learning.
\newblock {\em Machine Learning}, 28(1):41--75, 1997.

\bibitem{kirkpatrick2017overcoming}
J.~Kirkpatrick et~al.
\newblock Overcoming catastrophic forgetting in neural networks.
\newblock {\em Proceedings of the National Academy of Sciences}, 114(13):3521--3526, 2017.

\bibitem{riemer2019learning}
M.~Riemer et~al.
\newblock Learning to learn without forgetting by maximizing transfer and minimizing interference.
\newblock In {\em International Conference on Learning Representations (ICLR)}, 2019.

\bibitem{chaudhry2019efficient}
A.~Chaudhry et~al.
\newblock Efficient lifelong learning with A-GEM.
\newblock In {\em International Conference on Learning Representations (ICLR)}, 2019.

\bibitem{chaudhry2018efficient}
A.~Chaudhry et~al.
\newblock An empirical evaluation of continual learning.
\newblock In {\em arXiv preprint arXiv:1805.08375}, 2018.

\bibitem{nguyen2018variational}
C.~V. Nguyen et~al.
\newblock Variational continual learning.
\newblock In {\em International Conference on Learning Representations (ICLR)}, 2018.

\bibitem{french1999catastrophic}
R.~M. French.
\newblock Catastrophic forgetting in connectionist networks.
\newblock {\em Trends in Cognitive Sciences}, 3(4):128--135, 1999.

\bibitem{mccloskey1989catastrophic}
M.~McCloskey and N.~J. Cohen.
\newblock Catastrophic interference in connectionist networks: The sequential learning problem.
\newblock {\em Psychology of Learning and Motivation}, 24:109--165, 1989.

\bibitem{yu2020gradient}
T.~Yu et~al.
\newblock Gradient surgery for multi-task learning.
\newblock {\em Advances in Neural Information Processing Systems (NeurIPS)}, 33:5824--5836, 2020.

\bibitem{sener2018multi}
O.~Sener and V.~Koltun.
\newblock Multi-task learning as multi-objective optimization.
\newblock {\em Advances in Neural Information Processing Systems (NeurIPS)}, 31, 2018.

\bibitem{zenke2017continual}
F.~Zenke et~al.
\newblock Continual learning through synaptic intelligence.
\newblock In {\em International Conference on Machine Learning (ICML)}, pages 3987--3995, 2017.

\bibitem{aljundi2018memory}
R.~Aljundi et~al.
\newblock Memory aware synapses: Learning what (not) to forget.
\newblock In {\em Proceedings of the European Conference on Computer Vision (ECCV)}, pages 139--154, 2018.

\bibitem{lopez2017gradient}
D.~Lopez-Paz and M.~Ranzato.
\newblock Gradient episodic memory for continual learning.
\newblock {\em Advances in Neural Information Processing Systems (NeurIPS)}, 30, 2017.

\bibitem{rusu2016progressive}
A.~A. Rusu et~al.
\newblock Progressive neural networks.
\newblock {\em arXiv preprint arXiv:1606.04671}, 2016.

\bibitem{ruder2017overview}
S.~Ruder.
\newblock An overview of multi-task learning in deep neural networks.
\newblock {\em arXiv preprint arXiv:1706.05098}, 2017.

\end{thebibliography}

\end{document}
